% Gabarit pandoc des corrigés publiés (brut et guidé). Contenu sans marque ni nom de site.
% Compilation : lualatex (latexmk -lualatex).
% Le Markdown reste la source unique : ne pas éditer ce fichier .tex à la main.
\documentclass[11pt,a4paper]{article}

% --- Moteur, polices (Latin Modern) et langue ---------------------------------
\usepackage{iftex}
\ifLuaTeX\else\errmessage{Ce document se compile avec lualatex}\fi
\usepackage{amsmath,amssymb,mathtools}
\usepackage{fontspec}
% mathrm=sym : \mathrm prend ses glyphes dans la police mathématique (ℓ, μ, Ω), comme KaTeX.
\usepackage[math-style=TeX,bold-style=TeX,mathrm=sym]{unicode-math}
% Repli pour les glyphes absents de Latin Modern (indices et exposants Unicode,
% flèches, filets des schémas en texte, symboles) : DejaVu, puis Latin Modern Math
% (ℓ, lettres mathématiques Unicode des légendes).
\directlua{
  luaotfload.add_fallback("repliTexte", {"DejaVuSerif:mode=harf;", "DejaVuSans:mode=harf;", "file:latinmodern-math.otf:mode=harf;"})
  luaotfload.add_fallback("repliMono", {"DejaVuSansMono:mode=harf;", "DejaVuSans:mode=harf;", "file:latinmodern-math.otf:mode=harf;"})
}
\setmainfont{Latin Modern Roman}[RawFeature={fallback=repliTexte}]
\setsansfont{Latin Modern Sans}[RawFeature={fallback=repliTexte}]
\setmonofont{Latin Modern Mono}[RawFeature={fallback=repliMono}]
\setmathfont{Latin Modern Math}
\usepackage[french]{babel}
\frenchsetup{SmallCapsFigTabCaptions=false}

% --- Mise en page --------------------------------------------------------------
\usepackage[a4paper,margin=2.2cm,headheight=14pt,footskip=1.2cm]{geometry}
\usepackage{microtype}
\setlength{\emergencystretch}{4em}
\tolerance=1500
\setlength{\parindent}{0pt}
\setlength{\parskip}{0.55\baselineskip plus 2pt minus 1pt}
\usepackage{xcolor}
% Gris nommés : babel-french rend « ! » actif dans le document.
\definecolor{gris2}{gray}{0.98}
\definecolor{gris3}{gray}{0.97}
\definecolor{gris5}{gray}{0.95}
\definecolor{gris35}{gray}{0.65}
\definecolor{gris40}{gray}{0.60}
\definecolor{gris45}{gray}{0.55}
\definecolor{gris50}{gray}{0.50}
\definecolor{gris55}{gray}{0.45}
\definecolor{gris60}{gray}{0.40}
\definecolor{gris65}{gray}{0.35}
\definecolor{gris70}{gray}{0.30}
\definecolor{gris80}{gray}{0.20}
\definecolor{gris85}{gray}{0.15}
\usepackage{graphicx}
\usepackage{float}
\usepackage{caption}
\captionsetup{font=small,labelformat=empty,justification=centering}
\usepackage{booktabs,longtable,array,multirow,calc}
\usepackage{siunitx}
\sisetup{locale=FR}
\usepackage[version=4]{mhchem}
\usepackage{listings}
\usepackage{fancyhdr}
\usepackage{lastpage}
\usepackage[most]{tcolorbox}
\usepackage{enumitem}
\setlist{nosep,leftmargin=*,topsep=2pt}
\usepackage[hidelinks,unicode]{hyperref}
\usepackage{bookmark}
\hypersetup{pdftitle={Agro-Véto BCPST — Anglais (LVO) 2026},pdfauthor={Rédigé par IA},pdfcreator={LaTeX},pdfsubject={Corrigé guidé}}
\usepackage{xurl}

% --- Pandoc : listes, tableaux, figures ----------------------------------------
\providecommand{\passthrough}[1]{#1}
\providecommand{\tightlist}{\setlength{\itemsep}{0pt}\setlength{\parskip}{0pt}}

\makeatletter
\def\maxwidth{\ifdim\Gin@nat@width>\linewidth\linewidth\else\Gin@nat@width\fi}
\def\maxheight{\ifdim\Gin@nat@height>0.55\textheight 0.55\textheight\else\Gin@nat@height\fi}
\makeatother
\setkeys{Gin}{width=\maxwidth,height=\maxheight,keepaspectratio}
\setcounter{secnumdepth}{-\maxdimen}
\providecommand{\pandocbounded}[1]{#1}
\usepackage{longtable,booktabs,array}

% --- Code (listings) -----------------------------------------------------------
\lstset{
  basicstyle=\ttfamily\small,
  columns=fullflexible,
  keepspaces=true,
  breaklines=true,
  breakatwhitespace=false,
  postbreak=\mbox{\textcolor{gris50}{$\hookrightarrow$}\space},
  frame=single,
  rulecolor=\color{gris35},
  framesep=4pt,
  xleftmargin=5pt,xrightmargin=5pt,
  backgroundcolor=\color{gris3},
  keywordstyle=\bfseries,
  commentstyle=\color{gris65},
  stringstyle=\color{gris80},
  showstringspaces=false,
  tabsize=4,
  extendedchars=true,
  aboveskip=6pt,belowskip=6pt,
}

% --- Macros des corrigés (KaTeX ↔ LaTeX) ---------------------------------------
% Les corrigés sont écrits pour KaTeX : on définit ici ce que KaTeX connaît et
% que LaTeX n'a pas, sans jamais retoucher le Markdown.
\providecommand{\R}{\mathbb{R}}
\providecommand{\N}{\mathbb{N}}
\providecommand{\Z}{\mathbb{Z}}
\providecommand{\Q}{\mathbb{Q}}
\providecommand{\C}{\mathbb{C}}
\providecommand{\dd}{\mathrm{d}}
\providecommand{\llbracket}{\lBrack}
\providecommand{\rrbracket}{\rBrack}
\providecommand{\lBrace}{\lbrace}
\providecommand{\rBrace}{\rbrace}
\providecommand{\Box}{\square}
\providecommand{\Diamond}{\lozenge}
\providecommand{\varnothing}{\emptyset}
\providecommand{\argmax}{\operatorname*{arg\,max}}
\providecommand{\argmin}{\operatorname*{arg\,min}}
\renewcommand{\leq}{\leqslant}
\renewcommand{\geq}{\geqslant}

% Équation hors texte : réduite à la largeur de la ligne si elle déborde.
\newsavebox{\boiteequation}
\newcommand{\equationajustee}[1]{%
  \sbox{\boiteequation}{$\displaystyle #1$}%
  \ifdim\wd\boiteequation>\linewidth
    \[\resizebox{\linewidth}{!}{\usebox{\boiteequation}}\]%
  \else
    \[#1\]%
  \fi}

% --- Encadrés pédagogiques (sobres, lisibles en noir et blanc) -----------------
\tcbset{
  encadre/.style={
    enhanced,breakable,sharp corners,
    colback=white,colframe=gris55,
    boxrule=0.4pt,leftrule=2.2pt,
    left=6pt,right=6pt,top=3pt,bottom=3pt,
    before skip=6pt,after skip=6pt,
    fonttitle=\bfseries\small,coltitle=black,
    attach title to upper,after title={\quad},
    parbox=false,
  }
}
\newtcolorbox{indication}[1]{encadre,colframe=gris45,colback=gris2,title={Indication~#1}}
\newtcolorbox{methode}{encadre,colframe=gris70,title={Méthode}}
\newtcolorbox{piege}{encadre,colframe=gris85,boxrule=0.8pt,leftrule=2.2pt,title={Piège}}
\newtcolorbox{jury}{encadre,colframe=gris60,colback=gris5,title={Jury}}
\newtcolorbox{encadrenote}{encadre,colframe=gris40}

% --- Titres de section ---------------------------------------------------------
\usepackage{titlesec}
\titleformat{\section}{\Large\bfseries}{}{0pt}{}
\titleformat{\subsection}{\large\bfseries}{}{0pt}{}
\titleformat{\subsubsection}{\normalsize\bfseries}{}{0pt}{}
\titlespacing*{\section}{0pt}{14pt plus 4pt minus 2pt}{6pt}
\titlespacing*{\subsection}{0pt}{10pt plus 3pt minus 2pt}{4pt}
\titlespacing*{\subsubsection}{0pt}{8pt plus 2pt minus 2pt}{2pt}

% --- En-tête et pied de page ---------------------------------------------------
\pagestyle{fancy}
\fancyhf{}
\fancyhead[L]{\small\textsc{Corrigé guidé} · Agro-Véto BCPST}
\fancyhead[R]{\small Anglais 2026}
\fancyfoot[L]{\footnotesize Corrigé rédigé par IA, non officiel}
\fancyfoot[R]{\footnotesize \thepage\,/\,\pageref*{LastPage}}
\renewcommand{\headrulewidth}{0.4pt}
\renewcommand{\footrulewidth}{0pt}
\fancypagestyle{plain}{\fancyhf{}\fancyfoot[R]{\footnotesize \thepage\,/\,\pageref*{LastPage}}\fancyfoot[L]{\footnotesize Corrigé rédigé par IA, non officiel}\renewcommand{\headrulewidth}{0pt}}

\begin{document}
\thispagestyle{plain}
\begin{center}
  {\small\scshape Corrigé guidé}\par\smallskip
  {\LARGE\bfseries Agro-Véto BCPST — Anglais (LVO) 2026\par}\medskip
  {\large Agro-Véto · BCPST · Anglais · 2026\par}
\end{center}
\vspace{2pt}\hrule\vspace{8pt}

\emph{Corrigé rédigé par IA, non officiel. Signalez toute erreur.}

\hypertarget{pruxe9sentation-du-sujet}{%
\section{Présentation du sujet}\label{pruxe9sentation-du-sujet}}

\textbf{Format.} L'épreuve dure 2 heures, sans document ni appareil. Elle comprend deux parties rédigées en anglais :

\begin{longtable}[]{@{}
  >{\raggedright\arraybackslash}p{(\columnwidth - 6\tabcolsep) * \real{0.2145}}
  >{\raggedright\arraybackslash}p{(\columnwidth - 6\tabcolsep) * \real{0.4232}}
  >{\raggedright\arraybackslash}p{(\columnwidth - 6\tabcolsep) * \real{0.2106}}
  >{\raggedright\arraybackslash}p{(\columnwidth - 6\tabcolsep) * \real{0.1216}}@{}}
\toprule\noalign{}
\begin{minipage}[b]{\linewidth}\raggedright
Partie
\end{minipage} & \begin{minipage}[b]{\linewidth}\raggedright
Nature
\end{minipage} & \begin{minipage}[b]{\linewidth}\raggedright
Longueur
\end{minipage} & \begin{minipage}[b]{\linewidth}\raggedright
Barème
\end{minipage} \\
\midrule\noalign{}
\endhead
\bottomrule\noalign{}
\endlastfoot
I --- Compréhension de l'écrit & réponse à une question portant sur deux articles, l'un en français, l'autre en anglais, « in your own words » & 220 mots ± 10 \% (198 à 242) & 10 points \\
II --- Expression écrite & argumentation sur l'un des deux sujets, au choix & 220 mots ± 10 \% (198 à 242) & 10 points \\
\end{longtable}

Pour chaque partie, le nombre de mots est indiqué à la fin, et le numéro du sujet choisi en partie II est précisé : le sujet annonce des points de pénalité en cas de non-respect de ces consignes. Ce corrigé traite les deux sujets d'expression, soit trois questions : la compréhension (question 1), le sujet 1 (question 2), le sujet 2 (question 3).

\textbf{Les deux textes.}

\begin{longtable}[]{@{}
  >{\raggedright\arraybackslash}p{(\columnwidth - 4\tabcolsep) * \real{0.3059}}
  >{\raggedright\arraybackslash}p{(\columnwidth - 4\tabcolsep) * \real{0.1376}}
  >{\raggedright\arraybackslash}p{(\columnwidth - 4\tabcolsep) * \real{0.5265}}@{}}
\toprule\noalign{}
\begin{minipage}[b]{\linewidth}\raggedright
Texte
\end{minipage} & \begin{minipage}[b]{\linewidth}\raggedright
Nature
\end{minipage} & \begin{minipage}[b]{\linewidth}\raggedright
Contenu utile
\end{minipage} \\
\midrule\noalign{}
\endhead
\bottomrule\noalign{}
\endlastfoot
Mickaël Tassart, « Mercosur : à vous de voter contre ! », \emph{Courrier picard}, 20 décembre 2025 (en français) & éditorial, texte d'opinion & la mobilisation des agriculteurs obtient le report de la signature à janvier ; les « ponts d'or » de Bruxelles pour la viande et le sucre sud-américains ; l'Allemagne veut exporter ses voitures, son industrie chimique s'impatiente ; Giorgia Meloni soutiendrait l'accord ; appel aux consommateurs à « voter contre » en n'achetant pas les produits importés \\
Ronan McGreevy, « ``We're going to get hammered'': Irish beef farmers fear for margins as Mercosur deal goes through », \emph{The Irish Times}, 10 janvier 2026 (en anglais) & reportage, témoignages d'éleveurs & accord adopté par l'UE malgré le vote contre de cinq pays dont l'Irlande ; une bonne année 2025 (prix du bœuf + 40 \%) ; 99 000 tonnes de bœuf à droits réduits ; écart d'échelle (Brésil : 238 millions de têtes, Irlande : six millions) ; ni identification ni traçabilité, hormones et antibiotiques interdits en Europe ; « raised the white flag » sur le bien-être animal ; les morceaux haut de gamme visés ; échanges déséquilibrés (produits agricoles contre machines, chimie, pharmacie) \\
\end{longtable}

Les deux articles ont trois semaines d'écart : le premier est écrit pendant le report de la signature, le second après l'adoption de l'accord. Ils se complètent : le texte français est un plaidoyer adressé au lecteur, le texte irlandais donne les chiffres et la parole aux éleveurs.

\textbf{Les deux sujets d'expression.} Le sujet 1 prolonge directement la conclusion de l'éditorial (« c'est le consommateur qui aura le dernier mot »). Le sujet 2 est plus abstrait : il oppose deux notions à définir, \emph{economic growth} et \emph{rural livelihoods}.

\textbf{Ce qu'il ne faut pas laisser passer.} Respecter la fourchette de mots et indiquer un décompte exact. Répondre à la question posée, qui porte sur la \textbf{nature des inquiétudes}, sans la transformer en « ces craintes sont-elles fondées ? ». Citer les sources, et les citer correctement. En partie II, argumenter : les exemples doivent servir un raisonnement.

\textbf{Gestion du temps.} Une répartition raisonnable : 25 minutes de lecture et de relevé, 30 minutes de rédaction pour la partie I, 45 minutes pour la partie II (choix du sujet, plan au brouillon, rédaction), 15 minutes pour relire la langue et compter les mots. Les deux parties pèsent autant.

\textbf{Ce que dit le rapport 2026.} Le jury rappelle d'abord l'esprit de l'épreuve : « il faut surtout répondre clairement et intelligemment aux questions, dans une langue correcte et structurée » (p.~1). Les copies qui ont vu la différence de nature entre les deux textes ont été récompensées (p.~2). Le jury reproche aux copies de « redire tout ce que disent les textes » au lieu de sélectionner (p.~3) et de réinventer la question (p.~4). En expression, il rappelle ses quatre principes (p.~4-5) et donne une proposition de réponse pour chaque exercice (p.~11-12). Le rapport ne publie aucune statistique de notes.

\hypertarget{partie-i-compruxe9hension-de-luxe9crit}{%
\section{Partie I --- Compréhension de l'écrit}\label{partie-i-compruxe9hension-de-luxe9crit}}

Une seule question, sur deux textes de langue et de nature différentes. Le travail consiste à relever, à classer, puis à restituer en anglais les raisons de l'inquiétude des agriculteurs, en renvoyant à chaque fois à la source.

\hypertarget{question-1-why-are-farmers-in-france-and-ireland-worried-about-the-eu-mercosur-agreement}{%
\subsection{Question 1 --- Why are farmers in France and Ireland worried about the EU-Mercosur agreement?}\label{question-1-why-are-farmers-in-france-and-ireland-worried-about-the-eu-mercosur-agreement}}

\begin{indication}{1}

Lisez la question mot par mot : \emph{why} (des raisons), \emph{farmers} (les agriculteurs, pas les gouvernements ni les industriels), \emph{in France and Ireland} (les deux articles doivent être exploités), \emph{worried} (des craintes, pas des faits avérés).

\end{indication}

\begin{indication}{2}

Faites un relevé en deux colonnes, une par article, puis regroupez les éléments en trois familles : la concurrence économique, l'inégalité des normes, le sentiment d'abandon politique. Chaque famille donne un paragraphe.

\end{indication}

\begin{indication}{3}

Pour le texte français, traduisez l'idée et non les mots : les « ponts d'or réservés par Bruxelles » désignent des conditions très avantageuses accordées par la Commission européenne, pas des ponts. Pour le texte anglais, reformulez les citations (\emph{there is no tagging, there is no traceability}) au lieu de les recopier.

\end{indication}

\hypertarget{solution}{%
\subsubsection{Solution}\label{solution}}

Relevé utile, classé par famille :

\begin{longtable}[]{@{}
  >{\raggedright\arraybackslash}p{(\columnwidth - 4\tabcolsep) * \real{0.1876}}
  >{\raggedright\arraybackslash}p{(\columnwidth - 4\tabcolsep) * \real{0.3578}}
  >{\raggedright\arraybackslash}p{(\columnwidth - 4\tabcolsep) * \real{0.4246}}@{}}
\toprule\noalign{}
\begin{minipage}[b]{\linewidth}\raggedright
Famille
\end{minipage} & \begin{minipage}[b]{\linewidth}\raggedright
\emph{Courrier picard}
\end{minipage} & \begin{minipage}[b]{\linewidth}\raggedright
\emph{The Irish Times}
\end{minipage} \\
\midrule\noalign{}
\endhead
\bottomrule\noalign{}
\endlastfoot
Concurrence sur les prix, revenus & viande et sucre sud-américains favorisés par Bruxelles & 99 000 tonnes à droits réduits ; marges fragiles après une seule bonne année ; morceaux haut de gamme visés ; Brésil 238 millions de têtes contre six millions en Irlande \\
Normes, santé & --- & ni identification ni traçabilité ; hormones, antibiotiques interdits en Europe ; contrôles inexistants ; bien-être animal sacrifié \\
Politique & l'Allemagne veut vendre des voitures, la chimie allemande s'impatiente ; signature seulement reportée ; Meloni favorable & accord adopté malgré l'Irlande ; produits agricoles contre produits industriels ; surprise du vote italien \\
\end{longtable}

On écarte ce qui ne répond pas à la question : la citation de Coluche, le parcours personnel de John Curran, les témoignages sur la qualité de la viande irlandaise (ce sont des atouts, pas des craintes).

Farmers in France and Ireland fear the EU-Mercosur agreement for economic, health and political reasons, as shown by an opinion piece in the Courrier picard (20 December 2025) and a more factual report in The Irish Times (10 January 2026).

Their first concern is price competition. As the Courrier picard notes, the deal will let more South American meat and sugar into Europe. Irish breeders quoted by The Irish Times, who have just enjoyed their first decent margins in decades, fear that 99,000 tonnes of cheaper beef, mostly high-value cuts, will wipe out their fragile profits. The gap in scale is huge: Brazil has 238 million cattle, Ireland six million.

Farmers also denounce unfair competition. According to these Irish farmers, South American cattle are neither tagged nor traced and are given hormones and antibiotics banned in Europe, while controls there are ineffective. They demand equivalent standards and fear that lower-quality, possibly unsafe meat will reach European consumers.

Finally, both articles suggest that farmers feel sacrificed by their leaders. The deal mainly serves European industry, since Germany hopes to sell more cars and its chemical sector is impatient, while Mercosur countries mostly export farm produce. Despite tractor protests in France and in Brussels, the signature was only postponed in December, and the deal was approved in January although Ireland voted against it. The Courrier picard therefore turns to consumers instead.

\textbf{(229 words)}

\begin{methode}

Une phrase d'ouverture suffit : elle nomme les deux sources (titre du journal et date) et annonce les trois familles de raisons. Chaque paragraphe commence par une phrase qui en donne l'idée (\emph{Their first concern is price competition}), puis l'appuie sur un ou deux faits sourcés. Les connecteurs (\emph{also}, \emph{finally}, \emph{therefore}) rendent visible la hiérarchie.

\end{methode}

\begin{piege}

\emph{Beef} est indénombrable (\emph{99,000 tonnes of beef}, jamais \emph{beefs}), \emph{cattle} est un pluriel sans \emph{s} (\emph{238 million cattle}), \emph{cow} est dénombrable. Le jury signale aussi les confusions \emph{consumer} / \emph{customer}, \emph{concurrence} / \emph{competition} et les barbarismes \emph{concurrency}, \emph{productors}, \emph{agricultors} (p.~8).

\end{piege}

\begin{piege}

Ne présentez pas comme une inquiétude ce qui est un motif de fierté. Dire que le bœuf irlandais est traçable ne répond pas à la question ; dire que le bœuf sud-américain ne l'est pas, si. Le jury fait exactement cette distinction (p.~4).

\end{piege}

\begin{jury}

Rapport 2026 : « il ne s'agit pas de redire tout ce que disent les textes, mais de sélectionner puis restituer les informations pertinentes relevées dans les deux articles » (p.~3). Les copies qui ont vu dans l'article du \emph{Courrier picard} « un article d'opinion, contrastant avec l'article plus informatif de l'Irish Times, ont été valorisées » (p.~2). Les sources se citent dans une phrase correcte : \emph{The Irish Times states that\ldots{}} ou une parenthèse, mais « il n'est pas convenable d'employer un chiffre comme sujet d'une phrase » (p.~3).

\end{jury}

\begin{jury}

Rapport 2026 (p.~4) : « la question portait sur la nature des inquiétudes des agriculteurs, et non sur leur éventuelle justification ». Le jury reproche les introductions qui posent une autre question (\emph{Are farmers justified in being worried?}) et conclut : « il faut répondre à la question posée, et non la réinventer. » Sa proposition de réponse (p.~11) compte 216 mots et suit le même découpage : économie, concurrence déloyale, qualité et santé, manque de soutien politique.

\end{jury}

\begin{piege}

Attention aux dates : les deux articles ont trois semaines d'écart, pas un an. Le jury note que les candidats qui ne les ont pas lues avec soin ont cru à un an d'intervalle, « ce qui a pu gêner leur appréciation de la situation » (p.~2). En anglais : \emph{on 20 December 2025} ou \emph{on December 20, 2025}, \emph{in January 2026}, avec la majuscule au mois.

\end{piege}

\hypertarget{partie-ii-expression-uxe9crite}{%
\section{Partie II --- Expression écrite}\label{partie-ii-expression-uxe9crite}}

Le candidat choisit un seul des deux sujets. Le premier est le plus accessible, parce que l'éditorial du \emph{Courrier picard} fournit l'accroche ; le second exige de définir ses termes. Les deux devoirs ci-dessous suivent le même schéma : contexte et question, deux ou trois paragraphes argumentés, conclusion qui répond.

\hypertarget{question-2-sujet-1-do-you-think-consumers-should-be-responsible-for-protecting-local-agriculture-through-their-purchasing-choices-why-or-why-not}{%
\subsection{Question 2 --- Sujet 1 : Do you think consumers should be responsible for protecting local agriculture through their purchasing choices? Why or why not?}\label{question-2-sujet-1-do-you-think-consumers-should-be-responsible-for-protecting-local-agriculture-through-their-purchasing-choices-why-or-why-not}}

\begin{indication}{1}

Repérez les mots qui orientent la réponse : \emph{should} (une obligation morale, pas un constat), \emph{responsible for} (c'est à eux de le faire, pas « coupables »), \emph{local agriculture} (européenne ou française face aux importations, pas seulement la vente directe à la ferme), \emph{Why or why not?} (justifier).

\end{indication}

\begin{indication}{2}

Partez de l'éditorial du \emph{Courrier picard} : il dit précisément que le consommateur « aura le dernier mot ». C'est l'accroche la plus pertinente. Puis demandez-vous : si ce n'est pas au consommateur seul, à qui d'autre ?

\end{indication}

\begin{indication}{3}

Un plan en trois temps fonctionne bien : oui en partie (la demande oriente l'offre), mais pas seul (contraintes budgétaires, information imparfaite), donc une responsabilité partagée (État, cantines, distributeurs).

\end{indication}

\hypertarget{solution-1}{%
\subsubsection{Solution}\label{solution-1}}

On suit ce plan : une introduction qui part de l'actualité (les manifestations contre le Mercosur, l'appel de l'éditorial) et pose la question, une concession appuyée sur un exemple britannique, une objection appuyée sur l'inflation, une solution de partage des responsabilités illustrée par une loi française, une conclusion nette.

\textbf{Subject 1}

In recent months, French and Irish farmers have protested against the EU-Mercosur agreement, which will bring cheaper South American beef onto European shelves. Since governments seem unable to stop it, an editorial in the Courrier picard urges shoppers to ``vote'' against the deal with their wallets. Should consumers be held responsible for protecting local agriculture?

To some extent, they should, because demand shapes supply. When shoppers choose local produce, they keep money in rural areas and reward higher animal welfare and environmental standards. In the United Kingdom, the Red Tractor logo, which many supermarkets display, helps consumers identify food produced by British farms. Buying local is therefore a concrete way of supporting farmers.

However, making consumers solely responsible would be unfair. When food prices rose by almost 20 per cent in a year in Britain in 2023, many low-income families simply could not afford to pay more. Moreover, labels are often confusing, and nobody can check where every ingredient in a ready meal comes from.

Responsibility should therefore be shared. Governments can impose the same standards on imports as on European products, and public canteens can set an example: in France, the EGalim law requires school canteens to serve a minimum share of quality and organic food. Retailers, too, can give local products more visibility.

In short, consumers can help to protect local agriculture, but they cannot carry this burden alone.

\textbf{(231 words)}

Choix justifiés :

\begin{itemize}
\tightlist
\item
  La question finale de l'introduction reprend le sujet sous une forme directe et correcte (\emph{Should consumers be held responsible\ldots?}), sans en changer le sens.
\item
  Le logo Red Tractor est un exemple du monde anglophone, que le jury bonifie ; la loi EGalim montre ce que peut faire la commande publique.
\item
  L'inflation alimentaire britannique de 2023 est arrondie (« almost 20 per cent ») : un ordre de grandeur sûr vaut mieux qu'un chiffre précis mal retenu.
\end{itemize}

\begin{methode}

Pour éviter le « oui » simpliste comme le « non » paresseux, cherchez \textbf{qui d'autre} pourrait agir. Le rapport cite les distributeurs, « plus particulièrement les grandes surfaces (supermarkets, big box retailers) », et les pouvoirs publics, qui peuvent imposer des réglementations ou accorder des « subventions (subsidies) » (p.~6).

\end{methode}

\begin{piege}

\emph{Responsible} s'écrit avec un \emph{i} et se construit avec \emph{for} (\emph{responsible for protecting}). Le jury relève que le mot est « bien trop souvent mal orthographié » (p.~6) et juge inacceptables les erreurs sur les mots donnés dans le sujet (p.~8). Les circuits courts se disent \emph{short supply chains}, et non « short circuits » (p.~5).

\end{piege}

\begin{jury}

Rapport 2026 (p.~5) : ce sujet « a été traité par un très grand nombre de candidats », qui ont souvent préféré des accroches créatives à l'éditorial français. Le jury met en garde : « Attention, donc, aux juxtapositions absurdes, aux contextualisations trop éloignées ou aux accroches pseudo-philosophiques vides. » Il donne en exemple une bonne accroche qui part de la pression des accords commerciaux sur les agriculteurs européens.

\end{jury}

\begin{jury}

Rapport 2026 (p.~5) : « Il faut argumenter, les exemples seuls ne peuvent pas tenir lieu de réponse. » Et : « Mieux vaut choisir peu d'exemples mais bien expliquer en quoi ils appuient l'argumentation. » Enfin, « Les exemples ayant trait au monde anglophone seront bonifiés. »

\end{jury}

\begin{piege}

Les questions directes se construisent avec l'auxiliaire : \emph{Should consumers be held responsible?}, \emph{Are consumers responsible?} Le jury relève \emph{*Do consumers are\ldots?} et \emph{*Does consumers should be responsable\ldots?} (p.~9), ainsi que la formule lourde \emph{To what extent should consumers be responsible?} (p.~4).

\end{piege}

\hypertarget{question-3-sujet-2-is-balancing-economic-growth-with-the-protection-of-rural-livelihoods-possible}{%
\subsection{Question 3 --- Sujet 2 : Is balancing economic growth with the protection of rural livelihoods possible?}\label{question-3-sujet-2-is-balancing-economic-growth-with-the-protection-of-rural-livelihoods-possible}}

\begin{indication}{1}

Définissez d'abord les deux notions. \emph{Economic growth} : industrialisation, infrastructures, investissements, emplois. \emph{Rural livelihoods} : les moyens d'existence des campagnes (revenus agricoles, terres, artisanat, services, environnement). Le mot central de la question est \emph{balancing} : on demande si un équilibre est possible, pas si la croissance détruit les campagnes.

\end{indication}

\begin{indication}{2}

Commencez par montrer la tension, avec des exemples où la croissance a coûté cher aux ruraux. Puis nuancez : quels outils permettent de concilier les deux ?

\end{indication}

\begin{indication}{3}

La réponse finale doit être claire : « oui, à condition que\ldots{} » ou « non, parce que\ldots{} ». La condition doit être concrète : investissements publics, labels de qualité, partage des revenus.

\end{indication}

\hypertarget{solution-2}{%
\subsubsection{Solution}\label{solution-2}}

On suit ce plan : une introduction qui montre pourquoi la question se pose, une première partie sur la tension (le barrage des Trois-Gorges, l'Amazonie), une deuxième sur les outils de conciliation (indications géographiques, énergies renouvelables, tourisme), une troisième sur le rôle des politiques publiques (Irlande, PAC), une conclusion qui répond.

\textbf{Subject 2}

In many countries, governments build motorways, factories and dams to stimulate growth, but these projects often take land, jobs and public services away from the countryside. Is balancing economic growth with the protection of rural livelihoods possible?

At first sight, the two goals seem incompatible, as growth tends to concentrate wealth in cities and large companies. China's Three Gorges Dam, for instance, produces huge amounts of electricity but forced more than a million people to leave their homes. In Brazil, the expansion of cattle ranching and soybean farming boosts exports while destroying the forest on which indigenous communities depend.

Yet growth does not have to be achieved at the expense of rural areas. Quality labels allow farmers to earn more from their know-how: Scotch whisky and Stilton cheese are protected names, which keeps value and jobs in their regions. Renewable energy and eco-tourism can also create local income, provided that rural communities share the profits.

Such a balance, however, depends on political choices. Ireland's National Broadband Plan, which aims to connect remote homes to high-speed internet, shows that public investment can help rural people take part in the modern economy. Similarly, the European Union's Common Agricultural Policy funds rural development as well as farming.

To conclude, the balance is possible, but only if growth is designed to be sustainable and inclusive rather than driven by profitability alone.

\textbf{(227 words)}

Choix justifiés :

\begin{itemize}
\tightlist
\item
  L'introduction ne reformule pas la question : elle la contextualise, puis la pose telle quelle.
\item
  Les deux exemples de tension sont lointains (Chine, Brésil), les exemples de conciliation sont européens et anglophones (Scotch whisky, Stilton, plan irlandais) : le devoir montre que la réponse dépend des choix politiques.
\item
  Le chiffre des Trois-Gorges est arrondi par excès de prudence (« more than a million people »).
\end{itemize}

\begin{methode}

Un sujet construit sur deux notions abstraites se traite en trois temps : définir, opposer, concilier. Chaque paragraphe commence par sa thèse (\emph{At first sight, the two goals seem incompatible}; \emph{Yet growth does not have to\ldots{}}; \emph{Such a balance, however, depends on political choices}), puis la prouve par un exemple bref.

\end{methode}

\begin{piege}

\emph{Economic growth} n'est pas \emph{economic evolution}, et \emph{rural people lives} est fautif : on écrit \emph{rural people's lives} ou mieux \emph{rural livelihoods}. Le jury relève ces deux erreurs dans une introduction qui déformait la problématique (p.~6-7).

\end{piege}

\begin{jury}

Rapport 2026 (p.~6) : « Peu de candidats ont choisi cette question qui, dans l'ensemble, n'a pas été bien comprise. » Le jury attendait que l'on montre les tensions : « L'idée importante était de montrer que la croissance profite souvent davantage aux villes ou aux grandes entreprises. » Il accepte les deux conclusions, à condition qu'elles soient nettes, par exemple : « oui, l'équilibre est possible, mais seulement si la croissance économique est pensée de manière durable et inclusive ».

\end{jury}

\begin{piege}

Le registre doit rester écrit : \emph{want to}, \emph{going to}, \emph{much more}, et non \emph{wanna}, \emph{gonna}, \emph{way more}, que le jury relève dans des textes argumentatifs (p.~10). Évitez aussi les confusions de connecteurs signalées par le jury, comme \emph{moreover} et \emph{however} (p.~7).

\end{piege}

\hypertarget{bilan}{%
\section{Bilan}\label{bilan}}

\textbf{À retenir pour la partie I.}

\begin{itemize}
\tightlist
\item
  Lire la question mot par mot : ici, elle porte sur la nature des inquiétudes, pas sur leur bien-fondé.
\item
  Relever dans les deux textes, classer en deux à quatre familles, rédiger un paragraphe par famille avec une phrase-thèse.
\item
  Nommer les sources sans ambiguïté (titre du journal, nature de l'article, date) et y renvoyer à chaque information.
\item
  N'ajouter ni contexte extérieur ni jugement personnel : le jury pénalise les éléments « qui ne se trouvent nulle part dans les textes » (p.~4).
\end{itemize}

\textbf{À retenir pour la partie II.}

\begin{itemize}
\tightlist
\item
  Contextualiser en une ou deux phrases liées au sujet, puis poser la question telle quelle ou sous une forme directe et correcte.
\item
  Argumenter : une idée par paragraphe, un exemple bref et expliqué, de préférence du monde anglophone.
\item
  Conclure en répondant clairement.
\end{itemize}

\textbf{Le décompte des mots.} Il doit être indiqué et exact : « Tout décompte manifestement erroné, et a fortiori toute estimation peu sincère, est susceptible d'être pénalisé » (p.~10). Le jury déconseille aussi les marques de comptage trop visibles dans la copie. Viser 220 à 230 mots laisse une marge si les conventions de décompte diffèrent (traits d'union, nombres).

\textbf{Langue à revoir.} Majuscules des nationalités et des mois (\emph{French}, \emph{Irish}, \emph{January}) ; noms indénombrables (\emph{beef}, \emph{information}) ; accord du présent simple (\emph{they want}, \emph{the union has}) ; temps et aspect (prétérit et \emph{present perfect}, pas de forme progressive dans \emph{the article explains}) ; pronom \emph{they} dont le référent doit rester clair ; orthographe des mots du sujet (\emph{Mercosur}, \emph{responsible}).

\textbf{Situations qui se retrouvent dans d'autres sujets.} Restituer les arguments d'un camp à partir d'un éditorial et d'un reportage. Discuter la responsabilité individuelle face à un problème collectif (consommation, climat, santé). Concilier deux objectifs opposés (croissance et environnement, libre-échange et protection des producteurs).

\end{document}
